% Source: https://artofproblemsolving.com/community/c3h1763561_circles_and_bananas
% Stable identity: aops-topic-1763561-q01
Everything happens in the plane $\mathbf{P}_1$, which contains $\ell$; the plane $\mathbf{P}_2$ plays no role. Tangency puts the centres on the radii: $O_2$ is the midpoint of $O_1T_1$, so $O_1O_2=1$ in the direction of $T_1$, and $O_3$ lies beyond $T_2$ on ray $O_1T_2$ with $O_1O_3=3$.

The direction $\overrightarrow{T_1O_1}$ makes $30^\circ$ with $\ell$, so the direction $\overrightarrow{O_1O_2}$ also makes $30^\circ$ with $\ell$ (it is the opposite direction), and $\overrightarrow{O_1O_3}$ is perpendicular to $\ell$. Hence the angle $\theta=\angle O_2O_1O_3$ is $90^\circ-30^\circ=60^\circ$ or $90^\circ+30^\circ=120^\circ$, depending on which side of $\ell$ the point $O_3$ is taken. By the law of cosines,
\[
O_2O_3^2=1^2+3^2-2\cdot1\cdot3\cos\theta=10-6\cos\theta ,
\]
which is $7$ for $\theta=60^\circ$ and $13$ for $\theta=120^\circ$. The configuration with $\theta=60^\circ$ is shown:

\begin{asy}
usepackage("amsmath");
size(6cm);
defaultpen(fontsize(10pt));
pair O1=(0,2.6);
real c=Cos(30),sn=Sin(30);
pair O2=O1+(-c,sn),T1=O1+2*(-c,sn),O3=O1+(0,3),T2=O1+(0,2);
pair H=O1+(2.6/sn)*(c,-sn);
draw((-3.6,0)--(5.2,0),linewidth(0.8));label("$\ell$",(5.2,0),E);
draw(circle(O1,2),linewidth(0.9));draw(circle(O2,1),blue);draw(circle(O3,1),red);
draw(T1--H,dashed);draw(O1--T2,dashed);
draw(O2--O3,heavygreen+linewidth(1));
draw(arc(H,1.1,150,180));label("$30^\circ$",H+2.0*dir(165));
dot(O1);dot(O2);dot(O3);dot(T1);dot(T2);
label("$O_1$",O1,dir(15));label("$O_2$",O2,dir(250));label("$O_3$",O3,E);label("$T_1$",T1,W);label("$T_2$",T2+0.25*dir(-40),dir(-40));
label("$\sqrt7$",(-1.5,4.95),heavygreen);
label("\mbox{}",(min(currentpicture,user=true)+max(currentpicture,user=true))/2,black); // keeps the text after the figure black when LaTeX typesets the labels inline
\end{asy}

The minimum is $\sqrt7$, so $m=\boxed{7}$.
